“Verwijder die voyeur!”
“Ik beoog geen gezeur,

maar het is het beding
voor deze bevalling

gesteld door de landheer.
Hij eist dat deze keer

een kundig medicus
deze unieke klus

van dichtbij begeleidt.”
Dus hield de boze meid

haar onderkleed nog aan,
toen ze in bad moest gaan.

“Zusje, dat kan zo niet!”
De bevallende griet,

zonder expertise,
zag dat haar chemise

langzaam naar boven dreef
en daar dobberend bleef.

“En hoe voed je haar straks?"
Hij grijnsde achterbaks

toen hij ‘Moeder’ haar zus
opzij schoof voor zijn klus.

Tijdens zijn inspectie
misviel zijn positie.

Door het vrije uitzicht
verdween zijn evenwicht

en kukelde de schurk
dwars door haar onderjurk

in het verwarmde bad
tegen haar blote gat.

“Eruit en oprotten!”
Ze liet met zich spotten,

behalve deze keer.
“Ik waarschuw de landheer

over uw misdraging.”
Waarop de dokter ging.

“De profeet leerde mij
dat elk kind zondevrij

uit het vruchtwater zwemt.
Wat de profeet ontstemd,

is de boze invloed,
die dan zijn schade doet.

Door een kind te dopen
kunnen we vast hopen

het snel af te schermen.
Ik wil het beschermen

direct na de baring.
De doopgeboorte ging

bij mijn eerste, Emma,
uiteindelijk prima.

Ze behield haar gaafheid.
Ze landde in zachtheid

zonder de kille schok,
waar ik als kind van schrok.

Water voelt zuiverend.
Dus wees niet huiverend

bezorgde zus van mij!
Ik voel me veilig en vrij

omarmd door een vloeistof.
Een holistisch lusthof

zonder wilde mannen.
Prettig en ontspannen.

Ik voel de warmtebron
als een eeuwige zon.

In warme harmonie
met mijn kalme baby

door water verbonden,
blijft zij vrij van zonden.”

Emmy ontschoot gezwind.
Het bad met waterkind

werd toen niet weggegooid.
Wellicht “baadde” het ooit.

Water beviel dus Emmy
al na haar conceptie,

tijdens haar geboorte
en in de pretsoorten

die ze daarna ervoer.
Ze gaf geen ene moer

om haar sproetige huid.
Scheen de zon eens voluit,

dan slonk de verbranding
nabij een zandafgraving

in het zandwinningsmeer.
Ze ontdekte nog meer.

Zoals stiekem baden,
zonder kameraden,

in haar blozend blootje
om het risicootje

dat je soms werd betrapt.
Werd ze dan naakt gesnapt,

dan was Emmy grijpbaar.
Zo stookte dit gevaar

haar bevroren lijf op.
Werd het ijs warme sop.

“Water schenkt ons leven.
In Job staat geschreven

hoe het verdampt, regent
en op ons afstevent.
Het in de zee begint
en rivieren verbindt

in een waterkringloop.
Als een eeuwige doop.

Een oneindig proces
met een zinrijke les;

ook het universum
is als een pendulum,

waarbij een contractie
volgt op een expansie.

Want Heraclitus zei:
"Alles stroomt",  "Panta Rhei".

Water glijdt door de hand.
Het mijdt elke weerstand.

Het doet niets oneerlijk
en handelt natuurlijk.

Het voorziet de Tao
zonder druk of ego.

Het bereikt een diepgang
zonder werklust of dwang

en kan ons dienen voor
de ziel als metafoor.”

Aldus sprak de profeet,
die ‘Moeder’ vaak nadeed.

Ze wilde verzinken
om zich te verdrinken

in de baarmoederschoot,
die haar ooit warmte bood.

Ze zonk in gedachten
naar de zomernachten

toen ze nog jonger was
en zwom in een rietplas.

Beschermd door het duister
en door het gefluister,

zwom ze met vriendinnen,
aanstaande boerinnen,

water rond haar lichaam.
Dit vond ze aangenaam

als zachte kieteling
en zwoele prikkeling.

In de wintermaanden
stookten de blootstaanden

eerst het koele water
en goten die later

traag over een vriendin
en ving men het op in

een teil waarin men stond.
Het circuleerde rond.

Een gloed over de rug
en benen gleed terug

in de bak waar voeten
het stroompje begroeten.

Meer nog dan de douches,
die tussen flamoesjes,

de meiden opwonden,
was als ze lang stonden

naar elkaar te gluren,
de kou te verduren.
De jeugdige wijven
lieten zich graag drijven

op de toppen van vuur
en de dalen van guur.

De kou was nu feller.
Waardoor Emmy sneller

haar verkleumde wit hoofd,
door de ijsven verdoofd,

in de schoot moest leggen
en vaarwel moest zeggen.

Het was al bijzonder.
Wie dompelt nou onder

in een bevroren meer?
Ze deed dat keer op keer.

De kou had haar geleerd
hoe men zich dan verweerd

door de ademhaling.
Waardoor haar focus ging

van de ven naar de lucht,
die kwam bij elke zucht.

De kilte verduren,
want leren vraagt schuren.

Met koudetolerant,
als troef achter de hand,

weerstond ze telkens meer
het Angelsaksisch weer.

Ontwaakte een oerkracht
die haar voordelen bracht

te zijn in de natuur
met steeds langere duur.

Dit keer duurde haar bad
korter dan zij tijd had.

Als brak en baadboter
klom de rode koter

in haar ontblote gat
uit het ijzig veenbad.

Want de speelvrouw verhuist
uit haar show vroeg: “Wat ruist

er door het struikgewas?”
Wie denkt u dat het was?
